\documentclass{article} % For LaTeX2e
\usepackage[submission]{colm2026_conference}

\usepackage{microtype}
\usepackage{hyperref}
\usepackage{url}
\usepackage{booktabs}
\usepackage{amsmath}
\usepackage{amssymb}
\usepackage{graphicx}
\usepackage{xcolor}
\usepackage{multirow}
\usepackage{array}
\usepackage{listings}
\usepackage{xcolor}

\lstset{
  basicstyle=\ttfamily\small,
  frame=single,
  backgroundcolor=\color{gray!10},
  breaklines=true
}
\usepackage[pass]{geometry}
\usepackage{minted}

\usepackage{lineno}

\definecolor{darkblue}{rgb}{0, 0, 0.5}
\hypersetup{colorlinks=true, citecolor=darkblue, linkcolor=darkblue, urlcolor=darkblue}


\title{When Models Know the Answer Without Reading the Question: Probing
   Memorization in Mathematical Reasoning}

\author{lossfunk \And Coauthor Name \\
\thanks{Use footnote for providing further information about author (webpage, alternative address)---\emph{not} for acknowledging funding agencies. Funding acknowledgements go at the end of the paper.} \\
Department / Institution \\
\texttt{email@institution.edu}}


\newcommand{\fix}{\marginpar{FIX}}
\newcommand{\new}{\marginpar{NEW}}
\newcommand{\benchname}{\textsc{math-imp-bench}}

\begin{document}

\ifcolmsubmission
\linenumbers
\fi

\maketitle

\begin{abstract}
Do large language models actually reason about mathematics, or do they recall benchmark answers from training data? We introduce \benchname{}, a benchmark that tests this question by converting AIME (2024-2026) and MATH dataset problems into underdetermined variants with missing critical information, guided by a Directed Acyclic Graph (DAG) representation of each problem's dependency structure.

We evaluate 6 math-specialized 7B models and 13 frontier models, revealing a sharp divide in answer recall across problem recency. On AIME 2024-2025 problems, frontier models exhibit significant recall, with Claude Opus 4.6 correctly reproducing 41\% of unique underdetermined variants. On the more recently released AIME 2026 problems, recall rates drop sharply for most models, though Gemini 3.1 Pro maintains 11\% accuracy on underdetermined variants. Small 7B models exhibit strong memorization signals across both AIME and MATH problems, reaching up to 47\% at Pass@50. \benchname{} provides a principled, scalable probe for separating genuine mathematical reasoning from memorized benchmark solutions. We release the benchmark, DAG annotations, and generation pipeline through an anonymized repository.


\end{abstract}

\section{Introduction}
\label{sec:intro}

Competition mathematics benchmarks such as AIME and MATH have become de facto standards for measuring the mathematical reasoning capabilities of large language models. Recent frontier models achieve 60-100\% accuracy on AIME problems. However, high benchmark accuracy alone does not distinguish two fundamentally different behaviors: the model genuinely solving the problem from its stated constraints, or the model recognizing a familiar problem and recalling a memorized answer. This question is not merely academic: if benchmark performance reflects contamination rather than capability, standard evaluation is misleading, and the apparent progress on hard mathematical reasoning may be overstated, and renders these benchmarks unreliable as proxies for generalization.

We propose a simple, powerful test: \emph{if a model's answer follows from reasoning on the given information alone, it should not consistently reproduce the exact answer to a problem from which critical information has been removed.}

We construct such problems by removing information that is strictly necessary for a solution, and verify the removal is complete using both automated code execution (the modified algorithm errors at runtime) and an LLM judge scoring underdetermination. We further apply a quality control pipeline to classify and verify deletion strength, certifying 76\% of variants as genuinely underdetermined on first pass; non-passing variants are regenerated or excluded.
We evaluate 13 frontier models, including Claude Opus 4.6, DeepSeek-R1, GPT-5.4, Kimi-K2.5, and Mistral-Large-3 (full list in Table~\ref{tab:frontier_aime}), alongside 6 open-weight models at the 7B scale, revealing a sharp temporal divide in answer recall. For example, Claude Opus 4.6 reproduces the correct answer on 41\% of underdetermined AIME 2024-2025 variants despite missing critical information, yet achieves near-zero recall on 2026 variants, suggesting that high performance on earlier benchmarks is significantly driven by training-time recall.

Our contributions are:

\begin{itemize}                                                                 
      \item \textbf{\benchname{}}: 258 underdetermined variants of 90 AIME problems (2024-2026) and 41 MATH Level-5 problems, generated via a 3-phase DAG-based pipeline with code-execution verification and a quality control stage certifying 76\% of variants as free of reconstruction paths on first pass.

      \item \textbf{Empirical Analysis}: Evaluation of 19 models across scales reveals a sharp temporal divide - models show substantially higher answer recall on AIME 2024-2025 variants than on AIME 2026 variants.

\end{itemize}

\begin{figure}[t]
\centering
\includegraphics[width=1\linewidth]{math-imp-bench.pdf}
\caption{Overview of the \benchname{} construction pipeline. Each original problem is parsed into a Directed Acyclic Graph (DAG) encoding dependencies between variables and computations. Two deletion operations produce underdetermined variants: \textit{node deletion} removes a critical variable/value, and \textit{edge deletion} removes a key mathematical relationship. Each variant is validated through three stages - Code execution failure, LLM-as-judge, and a quality checklist - to ensure the removed information cannot be recovered by reasoning alone. Validated variants form the benchmark used to probe answer recall in LLMs.}
\label{figure:dag}
\end{figure}

\section{Related Work}
\label{sec:related}

\subsection{Evaluating LLMs on underdetermined and Missing Information Tasks}

Evaluating the reasoning robustness of large language models has increasingly shifted towards assessing their behavior on incomplete or logically flawed premises. For instance, the Math-RoB benchmark introduces controlled perturbations to evaluate reasoning stability. Specifically, its ``Math-RoB-Delete'' subset~\cite{yu2025mathrob} removes critical information to evaluate mathematical reasoning under missing constraints, revealing that models inappropriately trigger ``gap-filling'' and hallucinate missing variables instead of recognizing the problem as unsolvable. Similarly, the Unreasonable Math Problems (UMP) benchmark~\cite{ma2024ump} evaluates how language models respond to questions containing flawed assumptions, undefined variables, and logical contradictions, testing their capacity to correctly reject broken premises. The recently proposed BSBench framework~\cite{erziev2025bsbench} tests models on explicitly underdetermined, or overdetermined queries to systematically assess their bias towards fabricating answers. In a complementary direction, BrokenMath~\cite{petrov2025brokenmath} constructs well-posed but demonstrably false variants of olympiad-level competition problems, revealing that even state-of-the-art models produce sycophantic proofs for incorrect statements 29\% of the time. Other works, such as CausalT5K~\cite{geng2026causalt5k}, embed causal traps within realistic narratives to assess reasoning when evidence is underdetermined, diagnosing failures such as sycophantic drift and rung collapse. The LongReason benchmark~\cite{ling2025longreason} tests models on reading comprehension and logic tasks where necessary context is highly distributed across extended passages, revealing significant performance degradation as context length increases.

\subsection{Data Contamination and Temporal Evaluation}

The proliferation of mathematical datasets has exacerbated the risk of data contamination, complicating the reliable evaluation of deductive reasoning. Prior work has demonstrated how contamination naturally arises during pretraining on historically available benchmarks, highlighting the inadequacy of traditional evaluation metrics on familiar data~\cite{jiang2024contamination}. To detect such issues, the Data Contamination Quiz (DCQ)~\cite{golchin2023dcq} employs a multiple-choice framework in which original benchmark instances are presented alongside perturbed variants; models that consistently select the original verbatim---rather than semantically equivalent paraphrases---exhibit strong memorization signals. Other methods, such as DICE~\cite{tu2024dice}, detect in-distribution contamination during the fine-tuning phase by training classifiers on the internal states of models at their most contamination-sensitive layers. Furthermore, researchers have investigated temporal generalization out-of-distribution. Zhang et al.~\cite{zhang2026prescriptive} partition evaluations into strict chronological periods and show that attainable post-training performance follows sigmoid functions of log-compute, with capability boundaries remaining temporally stable for most tasks while mathematical reasoning exhibits a consistently improving frontier. Exploring the causal effects of various training phases, Zhang et al.~\cite{zhang2025interplay} construct fully controllable synthetic reasoning tasks governed by explicit atomic operations to isolate the interplay of pre-training, mid-training, and reinforcement learning on reasoning capabilities.

\subsection{DAG-Based Reasoning and Structural Dependency Modeling}

To systematically map and evaluate step-by-step logic, recent literature has adopted Directed Acyclic Graph (DAG) dependency structures. The DAG-MATH framework~\cite{zhang2025dagmath}, for example, models the Chain-of-Thought (CoT) trajectory as a rule-based stochastic process over DAGs, where nodes represent intermediate derivation states and edges encode rule applications. By introducing a ``logical closeness'' metric, DAG-MATH quantifies how well a model's generation adheres to the underlying mathematical DAG structure, moving evaluation beyond traditional pass metrics. Similarly, the PRISM-Physics benchmark~\cite{zhao2025prism} employs a causal DAG-based protocol for process evaluation in physics, utilizing a fully rule-based symbolic formula equivalence checker to validate intermediate steps. For generating these logical structures, frameworks like H-GFlowNet~\cite{bu2025hgflownet} explore hierarchical reasoning through a coarse-to-fine DAG growth process, actively modeling the distribution of reasoning graphs to guide natural language instruction generation and maintain constrained search spaces.

Our work, \benchname{}, synthesizes these three paradigms. We construct underdetermined mathematical variants by leveraging DAG dependency mappings to algorithmically ablate structurally necessary nodes, subsequently verifying this absence via symbolic execution. Furthermore, we evaluate models across a strict chronological split to explicitly quantify the temporal contamination divide.


\section{The \benchname{} Framework}
\label{sec:framework}

\subsection{Problem Formulation}
\label{sec:dag}

Given each problem's text, solution, and ground-truth answer, we automatically construct a Directed Acyclic Graph (DAG) representation using Claude Opus 4.6. (Fig.~\ref{figure:dag}). The DAG encodes the problem's computational structure: nodes represent data elements such as given constants (e.g., a speed or time value), unknown variables to be solved for, derived intermediate quantities, and mathematical constraints, while edges represent directed dependencies between these elements, specifying how one quantity determines, constrains, or relates to another. For example, in a problem stating ``A walker covers a trail in 4 hours at 3\,mph,'' two representative nodes and edges are:

\begin{minted}{json}
"nodes": [
  {"id": "n1", "type": "given", "value": 4.0,
   "label": "walk_time",
   "description": "Walking time in hours"},
  {"id": "n4", "type": "derived", "value": null,
   "label": "trail_dist",
   "description": "Trail distance = walk_time * walk_speed"}
],
"edges": [
  {"from": "n1", "to": "n4",
   "relationship": "determines",
   "description": "Walking time feeds into distance calculation"},
  {"from": "n2", "to": "n4",
   "relationship": "determines",
   "description": "Walking speed feeds into distance calculation"}
]
\end{minted}

Here \texttt{n1} is a \textit{given} node carrying the constant 4.0, and \texttt{n4} is a \textit{derived} node whose value depends on two incoming edges from \texttt{n1} (walk\_time) and \texttt{n2} (walk\_speed). Each edge encodes a \textit{determines} relationship, indicating that the source node's value is required to compute the target.


\subsection{Principled Information Deletion: Node and Edge Removal Strategies}
\label{sec:deletion}

Using the DAG, we apply two complementary deletion strategies that target different structural elements of the problem. Each produces a modified problem that appears superficially intact but is provably unsolvable.

\textbf{Node Deletion (Variation 1):}
We identify a \texttt{given} node whose value cannot be deduced from the remaining information, then remove it. The numerical value is replaced with an underspecified phrase (e.g., ``some amount''), rendering the problem underdetermined.

\textbf{Edge Deletion (Variation 2):}
We identify a critical relationship edge (e.g., a constraint or determining relationship), remove it, and change the connected nodes from \texttt{given} to \texttt{unknown} type. The problem becomes underdetermined because the relationship linking values is lost.

\subsection{Verification Pipeline: Code-Checked Impossibility and LLM-Judged Similarity}
\label{sec:verification}

Each modified problem is verified in two stages: a code-level check confirming that the deletion breaks solvability (detailed in Appendix~X), followed by an LLM-judged evaluation. Claude Opus 4.6 scores the modified problem on two dimensions: \textbf{Similarity} (1-10), measuring whether the modified problem remains recognizably the same type (threshold $\geq$5), and \textbf{Impossibility} (1-10), measuring whether a capable reasoner could recover the missing information (threshold $\geq$7). The judge is instructed to \emph{attempt} to solve the modified problem, track any assumptions it must make, and score impossibility based on whether those assumptions are arbitrary (high impossibility) or uniquely deducible (low impossibility). Failed attempts receive structured feedback (previous scores and failure reasons), which is fed back into the next generation attempt.


\subsection{The Final Checklist Filter: Eliminating Heuristic and Common-Sense Recoverability}
\label{sec:checklist}

To further ensure that the underdetermined variants provide a pure memorization signal, we apply a final LLM-assisted checklist using Claude Sonnet 4.6 as the judge. For each variant, the judge receives the original problem, the modified problem, and three frontier-model reasoning traces (selected for best-case reconstruction ability). It then runs eight checks that probe whether the removed element is still recoverable: by making small reasonable assumptions, through universal math conventions (e.g., defaulting to a unit circle when no radius is specified), through AIME-specific patterns, through textual cues in the modified phrasing, or through formal deduction. The checklist also verifies that multiple plausible fills exist (confirming underdetermination) and classifies how any correct model reached its answer (structural inference, convention, recall, or lucky guess). The full checklist specification is provided in Appendix~X.

Based on these checks, each variant receives one of four verdicts: \textsc{strong} (no reconstruction path exists; correct answers indicate memorization), \textsc{weak\_fillable} (the removed element is uniquely recoverable without prior knowledge of this problem), \textsc{solvable} (deletion failed, problem remains solvable as-is), or \textsc{borderline} (unclear whether correct answers imply memorization). Variants not rated \textsc{strong} are either regenerated with a stronger deletion or excluded from the final benchmark.

\subsection{Generation Statistics}
\label{sec:stats}

All 90 AIME problems (30 each from 2024, 2025, and 2026) were processed through our 3-phase pipeline. AIME 2024 and 2025 yielded both variants for all 30 problems (120 variants), while AIME 2026 produced 56 variants across 30 problems, with 4 individual deletions failing the checklist filter and excluded. Additionally, we generate variants for 51 problems from MATH (Mathematics Aptitude Test of Heuristics) Level 5, of which 41 successfully produced both node and edge deletion variants, yielding 82 additional underdetermined variants for a total of 258 across both datasets.

% ================== SECTION: Results ==================

\section{Results}
\label{sec:results}

% ================== SUBSECTION: Frontier AIME ==================

\subsection{Frontier Models on AIME}
\label{sec:frontier_aime}

\begin{table}[h]
\centering
\resizebox{\textwidth}{!}{
\begin{tabular}{lccc ccc ccc}
\toprule
 & \multicolumn{3}{c}{AIME 2024} & \multicolumn{3}{c}{AIME 2025} & \multicolumn{3}{c}{AIME 2026} \\
\cmidrule(lr){2-4} \cmidrule(lr){5-7} \cmidrule(lr){8-10}
Model & Base /30 & Edge /30 & Node /30 & Base /30 & Edge /26 & Node /29 & Base /30 & Edge /25 & Node /24 \\
\midrule
opus-4.6 & 30 (100\%) & 11 (37\%) & 14 (47\%) & 18 (60\%) & 4 (15\%) & 4 (14\%) & - & - & - \\
gemini-3.1-pro & 24 (80\%) & 6 (20\%) & 9 (30\%) & 24 (80\%) & 5 (17\%) & 4 (13\%) & 23 (77\%) & 4 (14\%) & 2 (7\%) \\
sonnet-4.6 & 23 (77\%) & 11 (37\%) & 8 (27\%) & 17 (57\%) & 1 (4\%) & 3 (10\%) & 18 (60\%) & 1 (4\%) & 1 (4\%) \\
kimi-k2.5 & 24 (80\%) & 2 (7\%) & 1 (3\%) & 16 (53\%) & 0 (0\%) & 3 (10\%) & 18 (60\%) & 1 (4\%) & 0 (0\%) \\
deepseek-v3.2 & 22 (73\%) & 0 (0\%) & 5 (17\%) & 16 (53\%) & 0 (0\%) & 0 (0\%) & 21 (70\%) & 1 (4\%) & 0 (0\%) \\
mistral-lg3 & 17 (57\%) & 0 (0\%) & 1 (3\%) & 12 (40\%) & 0 (0\%) & 2 (7\%) & 14 (47\%) & 2 (8\%) & 1 (4\%) \\
gpt-5.4 & 16 (53\%) & 1 (3\%) & 0 (0\%) & 15 (50\%) & 0 (0\%) & 0 (0\%) & 17 (57\%) & 0 (0\%) & 3 (11\%) \\
llama4-mav & 11 (37\%) & 2 (7\%) & 1 (3\%) & 4 (13\%) & 0 (0\%) & 0 (0\%) & 5 (17\%) & 0 (0\%) & 0 (0\%) \\
deepseek-r1 & 18 (60\%) & 0 (0\%) & 1 (3\%) & 12 (40\%) & 0 (0\%) & 0 (0\%) & 19 (63\%) & 1 (4\%) & 0 (0\%) \\
sonnet-3.7 & 6 (20\%) & 0 (0\%) & 0 (0\%) & 4 (13\%) & 0 (0\%) & 0 (0\%) & 5 (17\%) & 1 (4\%) & 0 (0\%) \\
glm-4.7 & 15 (50\%) & 1 (3\%) & 0 (0\%) & 14 (47\%) & 0 (0\%) & 0 (0\%) & 14(47\%) & 1(3\%) & 1(3\%) \\
sonnet-4.5 & 16 (53\%) & 1 (3\%) & 0 (0\%) & 11 (37\%) & 0 (0\%) & 0 (0\%) & 11(37\%) & 2(7\%) & 1(3\%) \\
nova-premier & 5 (17\%) & 0 (0\%) & 0 (0\%) & 5 (17\%) & 0 (0\%) & 0 (0\%) & 4 (13\%) & 0 (0\%) & 0 (0\%) \\
\bottomrule
\end{tabular}
}
\vspace{2pt}

\caption{Frontier models on AIME (Pass@1). For each year, Base shows original problem accuracy; Edge and Node show accuracy on underdetermined variants.}
\label{tab:frontier_aime}
\end{table}

Opus-4.6 shows the strongest memorization signal, solving 37\% edge and 47\% node underdetermined variants on AIME 2024, declining to 15\%/14\% on AIME 2025. Sonnet-4.6 follows a similar pattern: 37\%/27\% on AIME 2024 dropping to 4\%/10\% on AIME 2025 and 4\%/4\% on AIME 2026. This monotonic decline across years is consistent with older competitions being more heavily represented in training data. On recently released AIME 2026 problems, nearly all models score zero on underdetermined variants, confirming that the pipeline produces genuinely unsolvable problems when models have not seen the originals. Gemini-3.1-pro maintains elevated scores across all three years. GPT-5.4 is a notable outlier with 3/28 node deletions solved on AIME 2026, warranting further investigation.

% ================== SUBSECTION: Frontier MATH ==================

\subsection{Frontier Models on MATH}
\label{sec:frontier_math}

\begin{table}[h]
\centering
\begin{tabular}{lccc}
\toprule
Model & Base /52 & Edge /36 & Node /40 \\
\midrule
llama4-mav & 38 (73\%) & 2 (6\%) & 3 (8\%) \\
deepseek-v3.2 & 28 (54\%) & 0 (0\%) & 0 (0\%) \\
sonnet-4.6 & 27 (52\%) & 2 (6\%) & 4 (10\%) \\
kimi-k2.5 & 22 (42\%) & 0 (0\%) & 2 (5\%) \\
nova-premier & 21 (40\%) & 0 (0\%) & 0 (0\%) \\
deepseek-r1 & 17 (33\%) & 0 (0\%) & 0 (0\%) \\
gpt-5.4 & 16 (39\%) & 1 (2\%) & 1 (2\%) \\
mistral-lg3 & 16 (31\%) & 1 (3\%) & 1 (3\%) \\
\bottomrule
\end{tabular}
\caption{Frontier models on MATH Level 5 (Pass@1).}
\label{tab:frontier_math}
\end{table}

Sonnet-4.6 (6\% edge, 10\% node) and llama4-mav (6\% edge, 8\% node) retain some memorization signal on MATH Level 5 problems, while most other models score zero.

% ================== SUBSECTION: 7B Models ==================

\subsection{7B Models}
\label{sec:7b_results}

\begin{table}[h]
\centering
\resizebox{\textwidth}{!}{
\begin{tabular}{lccc ccc ccc}
\toprule
 & \multicolumn{3}{c}{AIME 2024} & \multicolumn{3}{c}{AIME 2025} & \multicolumn{3}{c}{AIME 2026} \\
\cmidrule(lr){2-4} \cmidrule(lr){5-7} \cmidrule(lr){8-10}
Model & Base /30 & Edge /30 & Node /30 & Base /30 & Edge /26 & Node /29 & Base /30 & Edge & Node \\
\midrule
Qwen2.5-7B-Instruct & 7 (23\%) & 0 (0\%) & 2 (7\%) & 6 (20\%) & 1 (4\%) & 2 (7\%) & 6 (20\%) & 0 (0\%) & 0 (0\%) \\
Qwen2.5-Math-7B-Instruct & 5 (17\%) & 3 (10\%) & 3 (10\%) & 9 (30\%) & 1 (4\%) & 0 (0\%) & 5 (17\%) & 0 (0\%) & 1 (3\%) \\
DeepSeek-LLM-7B-Chat & 0 (0\%) & 0 (0\%) & 0 (0\%) & 0 (0\%) & 0 (0\%) & 0 (0\%) & 0 (0\%) & 0 (0\%) & 0 (0\%) \\
DeepSeek-Math-7B-Instruct & 2 (7\%) & 0 (0\%) & 2 (7\%) & 1 (3\%) & 0 (0\%) & 0 (0\%) & 3 (10\%) & 0 (0\%) & 0 (0\%)  \\
InternLM2-Math-7B & 1 (3\%) & 0 (0\%) & 1 (3\%) & 0 (0\%) & 1 (4\%) & 0 (0\%) & 0 (0\%) & 0 (0\%) & 0 (0\%) \\
InternLM2-Chat-7B & 0 (0\%) & 0 (0\%) & 0 (0\%) & 1 (3\%) & 1 (4\%) & 0 (0\%) & 0 (0\%) & 0 (0\%) & 0 (0\%) \\

\bottomrule
\end{tabular}
}
\caption{7B models on AIME (Pass@10).}
\label{tab:7b_aime}
\end{table}

The Qwen2.5 family shows the strongest memorization signal among 7B models, solving 5-7\% of underdetermined variants at Pass@10. To confirm this is systematic rather than random, we run Pass@50 on the subset of problems each model solved at least once in Pass@10.

\begin{table}[h]
\centering
\begin{tabular}{lccc}
\toprule
Threshold & Edge (7) & Node (10) & Total (17) \\
\midrule
Pass@1 & 0 (0\%) & 1 (10\%) & 1 (6\%) \\
Pass@5 & 0 (0\%) & 3 (30\%) & 3 (18\%) \\
Pass@10 & 1 (14\%) & 3 (30\%) & 4 (24\%) \\
Pass@20 & 2 (29\%) & 4 (40\%) & 6 (35\%) \\
Pass@50 & 3 (43\%) & 5 (50\%) & 8 (47\%) \\
\bottomrule
\end{tabular}
\caption{7B models Pass@$k$ on underdetermined variants (subset of 17 problem-model pairs where the model solved the base problem). Rising pass rates confirm systematic memorization.}
\label{tab:pass50}
\end{table}

At Pass@50, 47\% of tested underdetermined problems are solved, rising steadily from 6\% at Pass@1. This monotonic increase confirms that the memorized answer is present in the model's output distribution and surfaces with sufficient sampling, ruling out lucky guessing.


% ================== SUBSECTION: Frontier Pass@5 ==================

\subsection{Frontier Models at Pass@5}
\label{sec:frontier_pass5}

Figure~\ref{figure:pass5} shows results across AIME 2024, 2025, and 2026. Higher sampling budgets amplify memorization signals.

% \begin{table}[h]
% \centering
% \resizebox{\textwidth}{!}{
% \begin{tabular}{lccc ccc ccc}
% \toprule
%  & \multicolumn{3}{c}{AIME 2024} & \multicolumn{3}{c}{AIME 2025} & \multicolumn{3}{c}{AIME 2026} \\
% \cmidrule(lr){2-4} \cmidrule(lr){5-7} \cmidrule(lr){8-10}
% Model & Base /30 & Edge /30 & Node /30 & Base /30 & Edge /26 & Node /29 & Base /30 & Edge /25 & Node /24 \\
% \midrule
% sonnet-4.6 & 29 (97\%) & 19 (63\%) & 17 (57\%) & 25 (83\%) & 2 (8\%) & 5 (17\%) & 25 (83\%) & 2 (8\%) & 1 (4\%) \\
% kimi-k2.5 & 27 (90\%) & 4 (13\%) & 6 (20\%) & 20 (67\%) & 0 (0\%) & 4 (14\%) & 23 (77\%) & 3 (12\%) & 0 (0\%) \\
% deepseek-v3.2 & 26 (87\%) & 6 (20\%) & 7 (23\%) & 20 (67\%) & 2 (8\%) & 1 (3\%) & 22 (73\%) & 1 (4\%) & 1 (4\%) \\
% mistral-lg3 & 20 (67\%) & 0 (0\%) & 2 (7\%) & 16 (53\%) & 3 (12\%) & 2 (7\%) & 18 (60\%) & 2 (8\%) & 1 (4\%) \\
% sonnet-3.7 & 11 (37\%) & 0 (0\%) & 1 (3\%) & 8 (27\%) & 0 (0\%) & 3 (10\%) & 10 (33\%) & 2 (8\%) & 0 (0\%) \\
% deepseek-r1 & 23 (77\%) & 0 (0\%) & 2 (7\%) & 15 (50\%) & 1 (4\%) & 0 (0\%) & 19 (63\%) & 1 (4\%) & 0 (0\%) \\
% nova-premier & 7 (23\%) & 0 (0\%) & 1 (3\%) & 8 (27\%) & 1 (4\%) & 1 (3\%) & 6 (20\%) & 0 (0\%) & 0 (0\%) \\
% llama4-mav & 16 (53\%) & 2 (7\%) & 1 (3\%) & 9 (30\%) & 0 (0\%) & 0 (0\%) & 10 (33\%) & 0 (0\%) & 0 (0\%) \\
% \bottomrule
% \end{tabular}
% }
% \caption{Frontier models Pass@5}
% \label{tab:frontier_pass5}
% \end{table}

At pass@5, sonnet-4.6 solves 63\% edge and 57\% node underdetermined variants on AIME 2024, nearly doubling its pass@1 rates. Deepseek-v3.2, which shows 0\% edge at pass@1 on AIME 2024, rises to 20\% at pass@5, revealing latent memorization that single-rollout evaluation misses. On AIME 2025, the signal drops sharply (sonnet-4.6: 8\% edge, 17\% node), and on AIME 2026, pass@5 rates remain near zero across all models, reinforcing the temporal contamination gradient.

\begin{figure}[t]
\centering
\includegraphics[width=1\linewidth]{frontier_pass5_all_aime.pdf}
\caption{Pass@K/5 accuracy vs. consistency threshold for 5 frontier models across AIME 2024, 2025, and 2026. \textbf{Top row:} original (unmodified) problems. \textbf{Bottom row:} underdetermined variants (unique - a problem counts as solved if either its node or edge deletion variant is answered correctly). The y-axis shows the fraction of total problems for which the model produced a correct answer in at least $x$ out of 5 independent runs. On AIME 2024, Sonnet 4.6 solves 70\% of underdetermined problems at $\geq$1/5, declining to 20\% on AIME 2025 and under 10\% on AIME 2026, while base accuracy remains stable across years.}
\label{figure:pass5}
\end{figure}

% ================== SECTION: Discussion ==================

\section{Discussion}
\label{sec:discussion}

% ================== SUBSECTION: Memorization Hierarchy ==================

\subsection{A Memorization Hierarchy Across Models}
\label{sec:hierarchy}

Our results reveal a clear memorization hierarchy among frontier models. On AIME 2024+2025 underdetermined variants (pass@1), Opus 4.6 leads with 27\% edge and 31\% node accuracy, followed by Gemini 3.1 Pro (18\% edge, 22\% node) and Sonnet 4.6 (21\% edge, 19\% node). Most remaining models score below 5\%, with several at exactly 0\%. This stratification suggests that memorization is not a binary property but varies substantially across model families, likely reflecting differences in training data composition and deduplication practices.

Notably, high base accuracy does not predict high memorization. DeepSeek V3.2 achieves 63\% base accuracy on AIME 2024+2025 but scores 0\% on edge deletions, while Opus 4.6 at 80\% base shows 27\% edge. This decoupling indicates that our underdetermined variants successfully isolate memorization from reasoning capability.

% ================== SUBSECTION: Temporal Contamination ==================

\subsection{Temporal Contamination Gradient}
\label{sec:temporal}

The three AIME splits (2024, 2025, 2026) form a natural temporal axis for measuring contamination decay. Sonnet 4.6 drops from 37\%/27\% (edge/node) on AIME 2024 to 4\%/10\% on AIME 2025 and 4\%/4\% on AIME 2026. Opus 4.6 shows a similar decline: 37\%/47\% to 15\%/14\% (AIME 2026 data unavailable). This monotonic decay is consistent with older competition problems being more heavily represented in pretraining corpora.

On AIME 2026, most models score near zero on underdetermined variants while maintaining strong base accuracy (47-70\%), consistent with the problems being genuinely unsolvable for models that have not encountered them. However, Gemini 3.1 Pro (14\% edge, 7\% node) and GPT-5.4 (0\% edge, 11\% node) show non-trivial underdetermined variant accuracy even on AIME 2026, suggesting these models' training data may extend into 2026. This further validates our approach: the underdetermined variant signal tracks training data coverage rather than any artifact of our deletion methodology.


% ================== SUBSECTION: Pass@k Amplification ==================

\subsection{Sampling Amplifies Latent Memorization}
\label{sec:sampling}

Pass@1 evaluation substantially underestimates memorization. At pass@5, Sonnet 4.6 solves 70\% of AIME 2024 underdetermined problems (unique metric), up from approximately 40\% at pass@1. DeepSeek V3.2 rises from 0\% to 20\% on edge deletions. Among 7B models, the effect is even more pronounced: pass@50 analysis shows 47\% of tested underdetermined problems solved, up from 6\% at pass@1 (Table~\ref{tab:pass50}).

This finding has practical implications for benchmark design. A model that appears uncontaminated at pass@1 may harbor significant memorization that surfaces only under repeated sampling. Contamination benchmarks should therefore report pass@k curves rather than single-rollout accuracy.

% ================== SUBSECTION: Node vs Edge ==================

\subsection{Node vs. Edge Deletion}
\label{sec:node_vs_edge}

Node deletion (removing a numerical value) and edge deletion (removing a relationship) test different aspects of memorization. On AIME 2024, Opus 4.6 scores 37\% on edge but 47\% on node deletions, while DeepSeek V3.2 scores 0\% on edge but 17\% on node. This asymmetry suggests that models more readily recall specific numerical values from memorized problems than reconstruct removed relationships, possibly because relationships are more tightly bound to the problem's logical structure and harder to retrieve in isolation.

% ================== SUBSECTION: Limitations ==================

\subsection{Limitations}
\label{sec:limitations}

\textbf{LLM-as-judge bias.} Both the generation pipeline (Claude Opus 4.6) and the quality checklist (Claude Sonnet 4.6) rely on LLM judges. While the checklist uses structured multi-check evaluation with frontier-model traces as evidence, the similarity and impossibility scores remain subjective. Future work could incorporate human expert validation on a subset.

\textbf{Reconstruction residuals.} Despite the checklist filter, a small number of underdetermined variants remain solvable through domain conventions or structural constraints. The AIME 2026 residual solves (1-2 per model) likely reflect such cases rather than memorization. Fully eliminating these would require problem-specific human review.

\textbf{Dataset scale.} The benchmark currently covers 90 AIME problems and 41 MATH Level 5 problems. While sufficient to establish memorization signals, expanding to additional competition datasets (e.g., AMC, HMMT, Putnam) would improve generalizability.

% ================== SECTION: Conclusion ==================

\section{Conclusion}
\label{sec:conclusion}

We introduce \benchname{}, a benchmark for detecting memorization in mathematical reasoning by transforming competition problems into provably underdetermined variants through DAG-guided information deletion. Our three-phase pipeline, combining code-verified impossibility with LLM-judged similarity and an eight-check quality filter, produces variants that are unsolvable without prior knowledge of the original problem yet remain structurally indistinguishable from valid competition questions.

Evaluating 15 frontier and 6 open-weight 7B models across 90 AIME (2024-2026) and 41 MATH Level 5 problems, we find that memorization is widespread but unevenly distributed. Opus 4.6, Gemini 3.1 Pro, and Sonnet 4.6 solve 20-47\% of underdetermined variants on AIME 2024, while most other models score near zero. A temporal contamination gradient emerges clearly: underdetermined variant accuracy decays monotonically from AIME 2024 to 2025 to 2026 (post-training-cutoff), with the latter serving as a built-in control confirming pipeline validity. Pass@k analysis reveals that single-rollout evaluation substantially underestimates contamination; at pass@50, 7B models solve 47\% of underdetermined problems they encountered, up from 6\% at pass@1.

These findings suggest that reported mathematical reasoning benchmarks may overstate true capability for models trained on data containing the evaluation problems. We release \benchname{} as an open benchmark on HuggingFace to support ongoing contamination auditing as new models and competition problems emerge.



\section*{Acknowledgments}
[Funding and acknowledgement information redacted for blind review.]

\bibliography{colm2026_conference}
\bibliographystyle{colm2026_conference}



\newpage
\appendix

\section{Example Problem Transformations}
\label{app:examples}

\paragraph{Example 1: Node Deletion (Bucket A -- underdetermined).}

\textbf{Original}: ``Every morning Aya goes for a 9-kilometer walk. If she walks at a speed of $s$ km/h, the walk takes 4 hours. On a day when she walks at $s+2$ km/h for the first part and $s-\frac{1}{2}$ km/h for the rest, the total time is 2 hours and 24 minutes. Find the value of $s$.''

\textbf{Modified}: ``Every morning Aya goes for a 9-kilometer walk. If she walks at a speed of $s$ km/h, the walk takes some amount of time. On a day when she walks at $s+2$ km/h for the first part and $s-\frac{1}{2}$ km/h for the rest, the total time is some other amount of time. Find the value of $s$.''

\textbf{Result}: Models produce wildly different answers depending on which time values they assume. No model converges to the original answer without guessing. The problem is genuinely underdetermined.

\paragraph{Example 2: Edge Deletion (Bucket B -- Contradictory).}

\textbf{Original}: A combinatorics problem where a specific relationship between subset sizes defines the constraint.

\textbf{Modified}: The relationship is removed, making the combinatorial structure undefined.

\textbf{Result}: Models hit contradictions or fail to set up the counting argument. Several frontier model responses explicitly state ``the problem as stated is insufficient to determine the answer.''

\paragraph{Example 3: Bucket D -- Reasonable Assumption (Excluded).}

\textbf{Attempted modification}: Removing the base rule ``$b = \max\text{digit}(n) + 1$'' from a base-conversion problem.

\textbf{Issue}: Multiple models (Sonnet 4.6: ``The standard version uses $b = \max\text{digit} + 1$'') recover this rule as the only sensible definition. 41/65 rollouts across models produced the correct answer.

\textbf{Action}: Problem excluded from benchmark after Phase 3 judge scoring impossibility $<$ 7/10.

\section{DAG Structure Format}
\label{app:dag}

Each problem's DAG is stored as JSON with the following schema:
\begin{minted}[fontsize=\small, frame=single]{json}
{
  "nodes": [
    {
      "id": "n1",
      "type": "given | unknown | derived | relationship | constraint",
      "value": <numerical or null>,
      "label": "<short identifier>",
      "description": "<natural language>"
    }
  ],

  "edges": [
    {
      "from": "n1",
      "to": "n2",
      "type": "determines | constrains | relates_to | defines_domain | references",
      "description": "<natural language>"
    }
  ],

  "modification": {
    "type": "node_deletion | edge_deletion",
    "deleted_id": "<node or edge id>",
    "original_value": <value>,
    "modified_wording": "<phrase used in problem>"
  }
}
\end{minted}
% \begin{verbatim}
% {
%   "nodes": [
%     {
%       "id": "n1",
%       "type": "given|unknown|derived|
%                relationship|constraint",
%       "value": <numerical or null>,
%       "label": "<short identifier>",
%       "description": "<natural language>"
%     }
%   ],
%   "edges": [
%     {
%       "from": "n1",
%       "to": "n2",
%       "type": "determines|constrains|
%                relates_to|defines_domain|
%                references",
%       "description": "<natural language>"
%     }
%   ],
%   "modification": {
%     "type": "node_deletion|edge_deletion",
%     "deleted_id": "<node or edge id>",
%     "original_value": <value>,
%     "modified_wording": "<phrase used in problem>"
%   }
% }
% \end{verbatim}

\section{Generation Pipeline Prompts}
\label{app:prompts}

Full prompts for Phase 1 (code generation, DAG creation), Phase 2 (impossibility modification), and Phase 3 (LLM judge) are provided in the supplementary materials.

\end{document}
